\section*{Chapitre \ref{ch:rc:credit-and-debit-valuation-adjustments} --- Ajustements de valeur pour risque de crédit et de risque de débit}

\begin{solution}{exo:rc:credit-and-debit-valuation-adjustments:1}
$0.6\times2\,000\,000\times(1-e^{-0.02\times5}) = \text{USD}~114\,195$.
\end{solution}

\begin{solution}{exo:rc:credit-and-debit-valuation-adjustments:2}
$6.44\text{ millions}\times0.0175\times7.45 = \text{USD}~840\,420$, contre 891\,342 en valeur exacte~: l'approximation
utilise une exposition non actualisée, moyennée dans le temps, et le spread à dix ans pour chaque date,
et elle fonctionne parce que le spread intègre déjà $\lambda(1-R)$.
\end{solution}

\begin{solution}{exo:rc:credit-and-debit-valuation-adjustments:3}
L'\omterm{def:rc:counterparty-exposure:netting}{ensemble de compensation} dérive en faveur de la contrepartie, si bien que la banque s'attend à devoir davantage qu'elle
ne doit recevoir~: sa DVA dépasse sa CVA, alors même que la contrepartie est plus risquée. La contrepartie voit
l'image inverse~: sa CVA sur la banque vaut USD~896\,328 en premier défaut, sa DVA 832\,416, soit
un coût net de 63\,912 qui constitue le bénéfice net de la banque.
\end{solution}

\begin{solution}{exo:rc:credit-and-debit-valuation-adjustments:4}
La CVA est linéaire en l'\omterm{def:rc:counterparty-exposure:profiles}{exposition positive attendue}, et la partie positive d'une somme est au plus
égale à la somme des parties positives~: la compensation permet aux valeurs des deux transactions de se neutraliser sur chaque trajectoire.
\end{solution}

\begin{solution}{exo:rc:credit-and-debit-valuation-adjustments:5}
$4\,254/(7.45\times10^{-4}) \approx \text{USD}~5.7$ millions de protection à dix ans. Restent non couvertes les
variations de l'exposition liées aux taux et au change (couvertes séparément), le \omterm{def:rc:rates-risk:crossgamma}{gamma croisé} entre crédit et
exposition, le \omterm{def:rc:counterparty-exposure:wwr}{risque de corrélation défavorable} (wrong-way risk), la base entre la courbe de la contrepartie et tout proxy, et
le \omterm{def:rc:reduced-form-credit:jtd}{risque de saut au défaut}~: une protection dimensionnée sur le \omterm{def:rc:reduced-form-credit:cs01}{CS01} ne compense pas la perte réelle en cas de défaut.
\end{solution}

\begin{solution}{exo:rc:credit-and-debit-valuation-adjustments:6}
Sous le CSA, les deux ajustements sont pilotés par les mouvements de la valeur sur la période de marge
en risque, qui sont symétriques~; sans CSA, la dérive de l'\omterm{def:rc:counterparty-exposure:netting}{ensemble de compensation} en faveur de la contrepartie
rendait importante l'exposition négative. Le CSA supprime davantage le côté le plus important.
\end{solution}

\begin{solution}{exo:rc:credit-and-debit-valuation-adjustments:7}
USD~766\,532, 1\,108\,459 et 1\,521\,524. Avec une dépendance très forte, les trajectoires où l'euro baisse
font presque sûrement défaut dès les premières années, quand l'exposition est encore faible, et la
survie moyenne doit être conservée, de sorte que les défauts sont tirés des expositions ultérieures, plus importantes~: la
CVA diminue après $b\approx7$.
\end{solution}

\begin{solution}{exo:rc:credit-and-debit-valuation-adjustments:8}
Un gain de DVA signifie que le crédit propre de la banque s'est dégradé~: le marché juge désormais plus vraisemblable qu'elle fasse défaut
sur ce qu'elle doit. Il s'agit d'une valeur au profit des créanciers, qui ne se réalise qu'en cas de défaut, et non d'un signe d'une meilleure
activité~; c'est pourquoi elle est déduite des fonds propres et généralement exclue du P\&L des traders.
\end{solution}

\begin{solution}{pb:rc:credit-and-debit-valuation-adjustments:1}
\textbf{1.} USD~1\,055\,437 avant, 1\,513\,008 après.
\textbf{2.} USD~457\,571, un gain de fair value des dérivés enregistré en compte de résultat.
\textbf{3.} Les engagements (l'exposition négative) sont désormais moins susceptibles d'être remboursés intégralement, si bien que
leur fair value baisse~: la banque gagne en tant que débitrice.
\textbf{4.} Environ USD~229 millions.
\textbf{5.} Les contreparties de la banque~: leur CVA sur la banque augmente du même montant.
\textbf{6.} Elle déduit la DVA des fonds propres de base de catégorie~1~: le gain ne compte pas en fonds propres.
\textbf{7.} À compter du 1 janvier 2018, après une montée en charge progressive à partir de 20~\% en 2014.
\textbf{8.} La CVA est une perte que réaliserait le défaut d'une contrepartie~; la DVA est un gain qui ne se réalise que
par le défaut propre de la banque, c'est-à-dire au moment où des fonds propres seraient nécessaires et font défaut.
\textbf{9.} Le ratio augmenterait à mesure que la banque s'affaiblit, l'inverse de ce que devraient faire les fonds propres.
\textbf{10.} Comme un élément non récurrent à retirer~: la plupart des analystes présentent les résultats avant les
ajustements liés au crédit propre.
\textbf{11.} Un swap de défaut sur soi-même ne paierait que si la banque a fait défaut et ne peut
pas encaisser~; personne ne le vendrait à un prix utile.
\textbf{12.} Vendre une protection sur des banques pairs, qui gagne lorsque les spreads bancaires se resserrent ensemble~;
elle ajoute le \omterm{def:rc:reduced-form-credit:jtd}{risque de saut au défaut} des pairs et la base entre les spreads de la banque et ceux des pairs.
\textbf{13.} La plupart des banques ne le font pas~: la couverture ajoute un risque réel pour protéger un chiffre comptable que
les fonds propres excluent déjà.
\textbf{14.} Sous le CSA à seuil nul, la DVA vaut USD~70\,490, si bien que le gain lié au même
élargissement est inférieur d'environ 93~\%.
\textbf{15.} Sa CVA sur la banque augmente du même USD~457\,571, une perte dans ses comptes.
\textbf{16.} Une valeur au profit des créanciers de la banque (qui supportent le déficit), et non de ses actionnaires.
\textbf{17.} Le présenter séparément et l'exclure de la performance du desk, ou l'imputer à une
fonction centrale~; les traders ne doivent pas être récompensés pour la dégradation du crédit de la banque.
\textbf{18.} La \omterm{def:rc:credit-and-debit-valuation-adjustments:bcva}{CVA bilatérale} de chaque partie est la DVA bilatérale de l'autre, si bien qu'un prix qui inclut
les deux est symétrique~; un prix unilatéral compte deux fois le même défaut, ou ne le compte pas du tout.
\textbf{19.} Résultat à retenir~: \emph{le gain lié à une prise de risque accrue}~: un élargissement de 50 points de base porte la
DVA de l'\omterm{def:rc:counterparty-exposure:netting}{ensemble de compensation} de USD~1\,055\,437 à 1\,513\,008, soit un gain de USD~457\,571.
\textbf{20.} Parce qu'elle croît précisément lorsque la capacité de la banque à absorber les pertes diminue.
\end{solution}

\begin{solution}{iq:rc:credit-and-debit-valuation-adjustments:1}
$\mathrm{CVA} = (1-R)\sum_k\mathrm{DEE}(t_k)\,\mathrm{PD}(t_{k-1},t_k)$~: \omterm{def:rc:reduced-form-credit:pd}{perte en cas de défaut}, \omterm{def:rc:counterparty-exposure:profiles}{exposition positive attendue} actualisée
de l'\omterm{def:rc:counterparty-exposure:netting}{ensemble de compensation} à chaque date, et la \omterm{def:rc:reduced-form-credit:pd}{probabilité de défaut} risque-neutre
de la contrepartie sur chaque période, issue de sa courbe de swaps de défaut~; indépendance supposée.
\iqlookfor{les trois termes et l'hypothèse d'indépendance.}
\end{solution}

\begin{solution}{iq:rc:credit-and-debit-valuation-adjustments:2}
La CVA de la contrepartie sur la banque~: la valeur de l'option de la banque de faire défaut sur ce qu'elle doit.
Elle est controversée parce qu'elle augmente lorsque le crédit de la banque se dégrade, ne peut se réaliser que par défaut,
ne se couvre pas directement et embellit les résultats en période difficile~; les régulateurs la déduisent des fonds propres.
\iqlookfor{la définition et les trois objections.}
\end{solution}

\begin{solution}{iq:rc:credit-and-debit-valuation-adjustments:3}
Couvrir les moteurs de l'exposition (delta et vega de taux) par des swaps et des swaptions, le crédit par
des swaps de défaut sur le nom ou un indice proxy dimensionné sur le \omterm{def:rc:reduced-form-credit:cs01}{CS01}, et gérer le \omterm{def:rc:reduced-form-credit:jtd}{risque de saut au défaut} par
une protection sur nom unique là où elle se négocie~; chercher un CSA ou des clauses de rupture~; surveiller le \omterm{def:rc:rates-risk:crossgamma}{gamma croisé}.
\iqlookfor{couvertures de marché et de crédit et leurs limites.}
\end{solution}

\begin{solution}{iq:rc:credit-and-debit-valuation-adjustments:4}
Simuler l'intensité de défaut de la contrepartie avec les facteurs de marché (une intensité stochastique corrélée
à ceux-ci, ou fonction de ceux-ci) et calculer des \omterm{def:rc:reduced-form-credit:pd}{probabilités de défaut} trajectoire par trajectoire~; ou décaler
les expositions conditionnellement au défaut~; ou recourir à des scénarios de stress. Conserver la courbe de survie moyenne
calibrée sur le marché.
\iqlookfor{simulation conjointe et calibration de la marginale.}
\end{solution}

\begin{solution}{iq:rc:credit-and-debit-valuation-adjustments:5}
La plupart des contreparties n'ont pas fait défaut, mais leurs spreads se sont fortement élargis et les expositions ont augmenté à mesure que
les taux et les changes ont évolué~: la CVA, valorisée au prix de marché, a progressé sur chaque \omterm{def:rc:counterparty-exposure:netting}{ensemble de compensation} non garanti à
la fois, alors que les défauts n'ont touché que quelques contreparties.
\iqlookfor{valorisation au prix de marché du crédit et défauts réalisés.}
\end{solution}

\begin{solution}{iq:rc:credit-and-debit-valuation-adjustments:6}
La clôture sans risque valorise les transactions résiliées sans le crédit de la partie survivante~; la clôture risquée
permet au prix de remplacement d'intégrer ce crédit. Avec une clôture risquée, la DVA de la partie survivante
demeure dans le \omterm{def:rc:credit-and-debit-valuation-adjustments:closeout}{montant de résiliation}, ce qui modifie les formules bilatérales et le recouvrement de
l'actif de la partie défaillante.
\iqlookfor{les deux conventions et leur effet sur la DVA.}
\end{solution}
